% Appendix C, typeset from project/README.md and info/examination.md.

\chapter{The group project}
\label{app:project}

This is the project's specification, collected in one place: the way of working, the structure, the
modules' interfaces, the test plan and the assessment. \Chapref{c10:ch} to \ref{c19:ch} are the
walkthrough; this is the contract, and where the two differ what is written here applies.

The task is to design, simulate and run a simplified CAN controller in VHDL, and make it reachable
from a microcontroller over SPI. The parallel class writes the C++ driver against the same registers,
so what is built is one half of a real hardware and software system, not an exercise ending in the
simulator.

When the project is finished the group has implemented a \textbf{bit timer} producing sample and
bit-end pulses, a \textbf{CRC-15 engine} that both computes and verifies the checksum, a
\textbf{transmit shift register} that serialises and inserts stuff bits, a \textbf{receive shift
register} that deserialises and removes them, a \textbf{CAN controller} whose state machine sequences
the frame's fields and handles arbitration, a \textbf{register bank} translating the controller's
pulses into pollable registers, an \textbf{SPI bridge} making the registers reachable from outside,
and a \textbf{top level} tying it all together into a node a microcontroller can talk to.

% ------------------------------------------------------------------------------------------
\section{Way of working}
\label{app:project:work}

The project is carried out in groups of \textbf{4--5 students}. The group is jointly responsible for
the whole project, and the way of working follows industrial practice.

\subsection{Git and GitHub}
\begin{itemize}
  \item Create a private repository and invite every group member and the supervisor as
    collaborators.
  \item Work in feature branches. One branch per module is a reasonable starting point.
  \item Commit often, with messages that say what the change does and not merely that it exists.
  \item \textbf{Never push directly to \code{main}.} All code reaches \code{main} through a pull
    request. Protect the branch in GitHub's settings so that the rule does not depend on everybody
    remembering it.
\end{itemize}

\subsection{Code review}
\begin{itemize}
  \item Every PR is reviewed by \textbf{at least one other group member} before it is merged.
  \item The reviewer checks: that the port order matches the interface table, that the module's
    testbench passes, that no values are hardcoded past \code{can_def}, and that every entity's
    interface and every important internal signal is commented.
  \item Put review comments in GitHub, not in chat or verbally. A review that does not exist in
    writing has not happened.
  \item A review that just says ``looks good'' is no review. Ask at least one real question, even
    when the code is good: \emph{why} it does what it does is more often interesting than
    \emph{that} it does it.
\end{itemize}

\subsection{Division of work}
Everybody is to write VHDL and everybody is to review VHDL. The division may be uneven in quantity,
not in kind. Document the division as you go in \file{CONTRIBUTORS.md} at the repository root.

% ------------------------------------------------------------------------------------------
\section{How the project is run}
\label{app:project:drive}

The project has \textbf{no milestones}. Every chapter from \chapref{c10:ch} onwards states a
\emph{target}, phrased as ``you should be roughly here''. The targets are not requirements and are
not graded.

That is a deliberate choice and not generosity. After \file{can_def.vhd}, the four leaf modules
\code{bit_timer}, \code{crc15}, \code{tx_shift_reg} and \code{rx_shift_reg} have \textbf{no
dependencies on each other}, and each has a testbench of its own handed out. A group of 4--5 can
therefore build them in parallel, in any order. Fixing an order would only get in the way of that.

But the dependency graph is not flat, and it is worth planning around:

\begin{textdrawing}
             meta_prev                                    (reads no package at all)
can_def  ──> bit_timer, crc15, tx_shift_reg, rx_shift_reg (the four leaves, parallel)
                                    └──> can_controller   (ch 16 transmit, ch 17 receive)
can_def  ──> register_bank                                (independent, buildable from ch 10)
spi_def  ──> spi_reg_bridge                               (independent, buildable from ch 10)
spi_slave                                                 (handed out, not written by the group)
    can_controller + register_bank + spi_reg_bridge + spi_slave ──> can_spi_node   (ch 19)
\end{textdrawing}

Two things follow from it:

\begin{itemize}
  \item \textbf{\code{can_controller} is the critical path.} It has no internal parallelism, it
    spans two sessions, and its testbench is a system testbench that either passes or does not
    \textemdash{} no partial green. Whoever is to own it should start reading \chapref{c16:ch}
    early, and the group should expect it to take the most time of anything.
  \item \textbf{\code{register_bank} and \code{spi_reg_bridge} are two full work streams that can
    start straight away.} One needs only \code{can_def}, the other only \code{spi_def}, and both
    have testbenches handed out. Nobody has to wait for the controller to have something to do, and
    a group that defers them to \chapref{c18:ch} and \ref{c19:ch} has put its tightest work last.
\end{itemize}

Two checkpoints are compulsory:

\begin{booktable}{lL}
  \thead{When} & \thead{What} \\
  \midrule
  \Chapref{c15:ch} & \textbf{Code review seminar.} Every group shows a merged PR and a review the
    group wrote, and talks about one thing in the way of working that has not worked. \\
  \Chapref{c19:ch} & \textbf{Final presentation and hand-in.} Green testbenches, a walkthrough of
    \code{can_spi_node}, bring-up on hardware as far as the group got, and the repository's
    history. \\
\end{booktable}

What is assessed is the end result and the way of working: the design, the Git history, the PRs and
the reviews. \textbf{Not the pace.} A group that is behind halfway through the project and hands in a
complete, well-reviewed node has done the project right.

% ------------------------------------------------------------------------------------------
\section{Project structure}
\label{app:project:layout}

Two flat directories, with the modules beside the testbenches that check them:

\begin{textdrawing}
controller/   can_def.vhd, meta_prev.vhd, bit_timer.vhd, crc15.vhd,
              tx_shift_reg.vhd, rx_shift_reg.vhd, can_controller.vhd
              + the six *_tb.vhd handed out
bridge/       register_bank.vhd, spi_reg_bridge.vhd, can_spi_node.vhd
              + spi_slave.vhd, spi_def.vhd, register_bank_tb.vhd,
                spi_reg_bridge_tb.vhd, handed out
ci/           build_project.sh, copied from the course repository
Makefile      with the build-project target
CONTRIBUTORS.md
.gitignore
\end{textdrawing}

The boundary between the two directories is the protocol: everything in \file{controller/} knows
about CAN and nothing about SPI, and \textbf{no module in either directory knows about both}.
\code{register_bank} is the middle: it knows neither the CAN frame format nor SPI but only register
semantics, and reads \code{can_def} purely for the port types.

The files handed out are fetched from the course repository and added unchanged. The group does not
write \file{spi_slave.vhd}: it is transport, not the course's subject, and it is read in
\chapref{c18:ch} instead.

% ------------------------------------------------------------------------------------------
\section{Naming conventions and coding style}
\label{app:project:style}

All the code and all the comments in the code are written in \textbf{English}. So are the
instructions.

\begin{itemize}
  \item Module and signal names in \code{snake_case}, for example \code{timer_enable}.
  \item Constants in capitals, for example \code{TICKS_PER_BIT}.
  \item \code{std_logic} and \code{std_logic_vector} for everything that is wires; \code{natural}
    for counters.
  \item No hardcoded values. Everything that is protocol lives in \code{can_def}.
  \item Comment every entity's interface and every important internal signal.
  \item Keep the interfaces stable. The testbenches are handed out and bind positionally.
\end{itemize}

\begin{warning}
\textbf{Positional port binding.} Every testbench handed out binds to its module
\textbf{positionally}, and so does every instantiation inside \code{can_controller}. The contract is
therefore the \textbf{order and the types} of the ports, exactly as the tables in
\secref{app:project:controller} and \ref{app:project:bridge} list them. The \emph{names} of the ports
are the group's own.

Swap two ports of the same type and nothing complains during analysis; the design simply misbehaves
in simulation. That is the price of letting the interface table be the single source of truth.
\end{warning}

% ------------------------------------------------------------------------------------------
\section{The shared package \code{can_def}}
\label{app:project:candef}

Written in \chapref{c10:ch}, in \file{controller/can_def.vhd}, and read by everything else.

\textbf{Constants:} \code{CLOCK_FREQ_HZ}, \code{BIT_RATE_HZ}, \code{TICKS_PER_BIT},
\code{SAMPLE_TICK}, \code{ID_WIDTH}, \code{DLC_WIDTH}, \code{DLC_MAX}, \code{CRC_WIDTH},
\code{DATA_WIDTH}, \code{MAX_RUN}, \code{CRC_POLY}.

\textbf{Subtypes:} \code{byte_t}, \code{id_t}, \code{crc_t}, \code{dlc_t}, \code{data_t}.

Target values for a DE0-CV with a 50 MHz system clock and CAN at 1 Mbit/s:

\begin{booktable}{llL}
  \thead{Constant} & \thead{Value} & \thead{Comment} \\
  \midrule
  \code{CLOCK_FREQ_HZ} & \code{50_000_000} & The system clock. \\
  \code{BIT_RATE_HZ} & \code{1_000_000} & 1 Mbit/s. \\
  \code{TICKS_PER_BIT} & \code{CLOCK_FREQ_HZ / BIT_RATE_HZ} & 50 clock cycles per CAN bit. Compute
    it, do not write it. \\
  \code{SAMPLE_TICK} & \code{TICKS_PER_BIT * 7 / 10} & The sample point, about 70 \% into the bit
    period. \\
  \code{ID_WIDTH} & \code{11} & Standard identifier. \\
  \code{DLC_WIDTH} & \code{4} & DLC is four bits. \\
  \code{DLC_MAX} & \code{8} & At most eight data bytes. \\
  \code{DATA_WIDTH} & \code{DLC_MAX * 8} & 64 bits of payload. \\
  \code{CRC_WIDTH} & \code{15} & CAN's CRC-15. \\
  \code{MAX_RUN} & \code{5} & The five rule for bit stuffing. \\
  \code{CRC_POLY} & \code{"100010110011001"} & Corresponds to $x^{15} + x^{14} + x^{10} + x^{8} +
    x^{7} + x^{4} + x^{3} + 1$. \\
\end{booktable}

That \code{TICKS_PER_BIT} and \code{SAMPLE_TICK} are \emph{computed} and not typed in matters: if
somebody changes the bit rate, everything should follow.

% ------------------------------------------------------------------------------------------
\section{The modules in \code{controller/}}
\label{app:project:controller}

The chapters describe the modules in detail; this is the port order collected, since that is what the
contract is.

\subsection*{\code{meta_prev} (\chapref{c12:ch})}
One generic, \code{WIDTH} (\code{natural := 1}), and three ports: \code{clock} (in), \code{async_in}
(in, \code{std_logic_vector(WIDTH-1 downto 0)}), \code{sync_out} (out, the same type). The
two-flip-flop synchroniser from \chapref{c4:ch}, written once with a width generic. It needs no
reset: after two clock cycles the contents are replaced anyway. Instantiated twice inside
\code{can_controller}, for \code{reset_n} and for \code{rx_bus}.

\subsection*{\code{bit_timer} (\chapref{c12:ch})}
\begin{longbooktable}{cll>{\raggedright\arraybackslash}p{0.24\linewidth}>{\raggedright\arraybackslash}p{0.30\linewidth}}{\thead{\#} & \thead{Port} & \thead{Direction} & \thead{Type} &
  \thead{Meaning}}
  1 & \code{clock} & in & \code{std_logic} & System clock. \\
  2 & \code{reset_s2_n} & in & \code{std_logic} & Synchronised, active-low reset. \\
  3 & \code{enable} & in & \code{std_logic} & Counts while \code{'1'}, holds while \code{'0'}. \\
  4 & \code{resync} & in & \code{std_logic} & Forces the counter back to tick 0. \\
  5 & \code{sample} & out & \code{std_logic} & Pulse at the sample point. \\
  6 & \code{bit_done} & out & \code{std_logic} & Pulse at the end of the bit period. \\
\end{longbooktable}

A counter from \code{0} to \code{TICKS_PER_BIT-1} pulsing \code{sample} at \code{SAMPLE_TICK} and
\code{bit_done} at the end of the period. \code{resync} is used once per frame, at SOF; in between
the timer free-runs.

\subsection*{\code{crc15} (\chapref{c13:ch})}
\begin{longbooktable}{cll>{\raggedright\arraybackslash}p{0.24\linewidth}>{\raggedright\arraybackslash}p{0.30\linewidth}}{\thead{\#} & \thead{Port} & \thead{Direction} & \thead{Type} &
  \thead{Meaning}}
  1 & \code{clock} & in & \code{std_logic} & System clock. \\
  2 & \code{reset_s2_n} & in & \code{std_logic} & Synchronised, active-low reset. \\
  3 & \code{clear} & in & \code{std_logic} & Clears the register without a reset. \\
  4 & \code{enable} & in & \code{std_logic} & Integrates \code{data} at this edge. \\
  5 & \code{data} & in & \code{std_logic} & The current bit. \\
  6 & \code{crc} & out & \code{crc_t} & The register's current value. \\
  7 & \code{valid} & out & \code{std_logic} & \code{'1'} when the register is all zeros. \\
\end{longbooktable}

A 15-bit register with feedback. The same engine both generates and checks: the transmitter feeds in
its bits and reads out \code{crc}; the receiver feeds in the same bits \emph{plus} the received CRC,
and \code{valid} goes high if everything matches. No mode switch is needed. \code{clear} is needed all
the same: an aborted frame leaves rubbish in the register.

\subsection*{\code{tx_shift_reg} (\chapref{c14:ch})}
\begin{longbooktable}{cll>{\raggedright\arraybackslash}p{0.24\linewidth}>{\raggedright\arraybackslash}p{0.30\linewidth}}{\thead{\#} & \thead{Port} & \thead{Direction} & \thead{Type} &
  \thead{Meaning}}
  1 & \code{clock} & in & \code{std_logic} & System clock. \\
  2 & \code{reset_s2_n} & in & \code{std_logic} & Synchronised, active-low reset. \\
  3 & \code{load} & in & \code{std_logic} & Loads \code{data}/\code{bit_count} \textbf{and} puts out
    the first bit the same cycle. \\
  4 & \code{data} & in & \code{byte_t} & The group, with its valid bits at the top. \\
  5 & \code{bit_count} & in & \code{std_logic_vector(3 downto 0)} & Number of valid bits, 1--8. \\
  6 & \code{shift} & in & \code{std_logic} & One pulse per CAN bit, from
    \code{bit_timer.bit_done}. \\
  7 & \code{tx_bit} & out & \code{std_logic} & The bit to drive next bit period. \\
  8 & \code{stuff} & out & \code{std_logic} & \code{'1'} on the very shift that puts out a stuff
    bit. \\
  9 & \code{done} & out & \code{std_logic} & Pulses when the group has been fully sent. \\
  10 & \code{bit_valid} & out & \code{std_logic} & \code{'1'} the cycle a new bit was put out. \\
\end{longbooktable}

Shifts out MSB first, counts identical bits in a row, and inserts a stuff bit of the opposite value
after \code{MAX_RUN} identical ones. \textbf{The stuffing state survives reloads}: the rule applies to
the bit stream, not to the individual group. \code{stuff} and \code{bit_valid} exist for
\code{can_controller}'s sake: the CRC is fed on \code{bit_valid and not stuff}, since a stuff bit is
not part of the checksum.

\subsection*{\code{rx_shift_reg} (\chapref{c15:ch})}
\begin{longbooktable}{cll>{\raggedright\arraybackslash}p{0.24\linewidth}>{\raggedright\arraybackslash}p{0.30\linewidth}}{\thead{\#} & \thead{Port} & \thead{Direction} & \thead{Type} &
  \thead{Meaning}}
  1 & \code{clock} & in & \code{std_logic} & System clock. \\
  2 & \code{reset_s2_n} & in & \code{std_logic} & Synchronised, active-low reset. \\
  3 & \code{sample} & in & \code{std_logic} & One pulse per CAN bit, from
    \code{bit_timer.sample}. \\
  4 & \code{rx_bus} & in & \code{std_logic} & The bus level, valid at the sample point. \\
  5 & \code{bit_count} & in & \code{std_logic_vector(3 downto 0)} & Number of \emph{real} bits to
    collect, 1--8. \\
  6 & \code{enable} & in & \code{std_logic} & Decides whether \code{sample} pulses are acted on. \\
  7 & \code{data} & out & \code{byte_t} & The collected group, right-justified. \\
  8 & \code{valid} & out & \code{std_logic} & Pulses when \code{bit_count} real bits have been
    collected. \\
  9 & \code{stuff_error} & out & \code{std_logic} & Pulses if an expected stuff bit does not
    invert. \\
  10 & \code{real_bit} & out & \code{std_logic} & The destuffed bit that was accepted. \\
  11 & \code{real_bit_valid} & out & \code{std_logic} & \code{'1'} when a real bit was accepted,
    \code{'0'} for a discarded stuff bit. \\
\end{longbooktable}

The transmitter's mirror image. After \code{MAX_RUN} identical bits in a row the next bit \emph{must}
be the opposite; if it is, it is discarded and the count starts over, and if it is not, that is a
\code{stuff_error}.

\subsection*{\code{can_controller} (\chapref{c11:ch}, \ref{c16:ch}, \ref{c17:ch})}
The top level of the CAN half, and the only block that knows the frame format. Fifteen ports, in
exactly this order, every input first:

\begin{textdrawing}
 1 clock     4 tx_id    7 rx_bus    10 bus_en   13 rx_data
 2 reset_n   5 tx_dlc   8 tx_done   11 rx_id    14 rx_valid
 3 tx_req    6 tx_data  9 tx_bus    12 rx_dlc   15 error
\end{textdrawing}

\textbf{The frame sequence} a transmission goes through: SOF (1 bit), the arbitration field (11-bit
identifier + RTR), the control field (IDE, r0, DLC = 6 bits), the data field (0--64 bits according to
DLC), CRC (15 bits + delimiter), ACK (slot + delimiter) and EOF (7 bits). Every field gets a load
state of its own before its shift state.

\textbf{The bus} is modelled as open drain: \code{bus_en} decides whether the node drives at all,
\code{tx_bus} what it drives. The bus value is \code{'1'} only if no node pulls it down. During the
ACK slot the transmitting node releases the bus, and a receiver with a valid CRC pulls it down.
\textbf{The transmitter never reads the ACK slot back}; a frame nobody acknowledges completes exactly
like an acknowledged one, and \code{tx_done} pulses all the same.

\textbf{Arbitration:} during the arbitration field every transmitted bit is compared against what is
actually on the bus. If the node sends a recessive bit and reads back a dominant one it has lost: it
sets \code{error}, pulses \code{tx_done}, leaves the transmitter role and goes idle. It does
\textbf{not} receive the frame that won.

That \code{tx_done} pulses here too is deliberate: the pulse means ``the transmission attempt is over
and the port is yours again'', not ``the frame arrived''. \code{error} says which. Without that pulse
a caller has no edge to wait for, and the register bank never gets its TX-ready bit back.

\textbf{\code{error}} goes high for four reasons: lost arbitration, \code{stuff_error}, a failed CRC
check, and a received remote frame. It stays high until the next transmission or reception starts.
\textbf{\code{tx_done}} and \textbf{\code{rx_valid}} are pulses of exactly one clock cycle, which is
precisely what makes the register bank necessary.

\subsection{Known limitations the driver notices}
\label{app:project:limits}

The simplifications against real CAN are listed in full in \secref{c19:sec:gaps}, and from the
protocol side in \canref{7}. Three of them show all the way out in the register map and therefore have
to be stated here too, since the parallel class writes code against them \textemdash{} they are the
same three \canbook{} highlights, and both classes read that book:

\begin{itemize}
  \item \textbf{A lost arbitration cannot be told from a successful transmission on STATUS bit 0.}
    Both give the bit back, since the controller pulses \code{tx_done} when it aborts as well. That is
    deliberate: the bit means ``the port is free'', not ``the frame arrived''. A driver that wants to
    know the outcome reads \code{ERROR_FLAGS} after bit 0 has come back.

    This was once a real bug, and it is worth knowing why it is not one any more. If bit 0 were set
    only by \code{tx_done}, and \code{tx_done} pulsed only at EOF, the bit would be left low after
    every lost arbitration and the node could not transmit again until it was cleared by the reset
    button \textemdash{} on a two-node bus, that is, in normal operation. Three things keep it from
    arising again, and they overlap deliberately: the controller pulses \code{tx_done} at aborts too
    (\chapref{c17:ch}), the bank sets bit 0 on a rising edge of \code{error} while the bit is low as
    well (\chapref{c18:ch}), and \code{TX_ABORT} exists as an escape hatch.
    \code{can_controller_tb} case 2 and \code{register_bank_tb} cases 10--11 check the three.
  \item \textbf{\code{error} is one bit for four causes.} The driver can see \emph{that} something
    went wrong, never \emph{what}. Retransmission logic distinguishing a lost arbitration from a
    broken CRC cannot be written against this interface.
  \item \textbf{No buffering on the receive side.} The controller holds one frame. If the next
    arrives before the previous has been read out, it is overwritten. The register bank passes the
    limitation on rather than hiding it; see appendix~\ref{app:protocol}.
\end{itemize}

% ------------------------------------------------------------------------------------------
\section{The modules in \code{bridge/}}
\label{app:project:bridge}

\subsection*{\code{register_bank} (\chapref{c18:ch})}
Sixteen ports, in this order:

\begin{textdrawing}
 1 clock        5 reg_wdata    9 rx_data     13 tx_id
 2 reset_s2_n   6 tx_done     10 rx_valid    14 tx_dlc
 3 reg_index    7 rx_id       11 error       15 tx_data
 4 reg_write    8 rx_dlc      12 reg_rdata   16 tx_req
\end{textdrawing}

Translates the controller's one-cycle pulses into sticky, pollable levels, and register writes into
events. The semantics in full are in appendix~\ref{app:protocol}; it is the contract with the parallel
class and must not be changed unilaterally. \code{reg_rdata} is \textbf{combinational}: latching a
read value is the bridge's responsibility, not the bank's.

\subsection*{\code{spi_reg_bridge} (\chapref{c19:ch})}
Ten ports, in this order: \code{clock}, \code{reset_s2_n}, \code{ss_active}, \code{rx_data},
\code{rx_valid}, \code{reg_rdata}, \code{tx_data}, \code{reg_addr}, \code{reg_wdata},
\code{reg_write}.

Implements the transaction from appendix~\ref{app:protocol}: one command byte plus four data bytes,
always exactly five. On a read the \textbf{register value is latched once} at the end of the command
byte; on a write nothing is committed until the fifth byte is complete. If \code{ss_active} goes low
before that, the transaction is abandoned entirely, with no side effects. \code{reg_addr}
\textbf{holds} the decoded index until the next command byte decodes a new one.

\subsection*{\code{can_spi_node} (\chapref{c19:ch})}
Nine ports, in this order:

\begin{textdrawing}
 1 clock      4 mosi     7 rx_bus
 2 reset_n    5 ss       8 tx_bus
 3 sclk       6 miso     9 bus_en
\end{textdrawing}

A purely structural top level. Instantiates \code{spi_slave}, \code{spi_reg_bridge},
\code{register_bank} and \code{can_controller}, plus a \code{meta_prev} for the reset. No functional
logic here.

The order is SPI's four wires as a group, \code{sclk}/\code{mosi}/\code{ss}/\code{miso}, and then the
CAN bus's three. \code{miso} therefore sits among the SPI ports and not with the other outputs, so
that the four pins stand in the same order as in the wiring table. That matters:
\code{can_spi_node_tb} binds positionally like everything else, and \code{miso} and \code{rx_bus} have
the same type, so swapping them passes analysis without a word and gives a node that never answers.

% ------------------------------------------------------------------------------------------
\section{Test plan}
\label{app:project:tests}

The group writes no testbenches in this project. Nine are handed out, and they are the contract: seven
module testbenches and two system testbenches.

\begin{booktable}{lL}
  \thead{Testbench} & \thead{Checks} \\
  \midrule
  \code{meta_prev_tb} & Two flip-flops of delay, for both one bit and a vector. \\
  \code{bit_timer_tb} & \code{sample} and \code{bit_done} at the right tick; \code{resync};
    \code{enable}. \\
  \code{crc15_tb} & CRC values for known sequences; \code{valid}; \code{clear}. \\
  \code{tx_shift_reg_tb} & MSB-first shifting, stuffing after five identical bits, the \code{done}
    timing. \\
  \code{rx_shift_reg_tb} & Deserialisation, destuffing, \code{stuff_error}, \code{real_bit}. \\
  \code{register_bank_tb} & Eleven cases, from the reset state to \code{TX_ABORT} and the two routes
    back to TX ready. \\
  \code{spi_reg_bridge_tb} & Four cases: a write, a read byte by byte, an abort on \code{ss_active},
    and a command byte with a reserved bit set. \\
\end{booktable}

\code{can_controller_tb} puts two nodes on a shared wired-AND bus and checks that one receives the
other's frame, and that arbitration is resolved correctly when both transmit at once. It is skipped by
the build until the controller has a receive path (appendix~\ref{app:sim}).

\code{can_spi_node_tb} does the same thing one layer out: two complete nodes on the same bus, but
every node driven by a modelled SPI master on its four pins. It plays back the protocol
specification's worked example through the whole chain and follows a frame from an SPI write on one
node to an SPI read on the other. Seven cases: the reset state, the worked example, invalid command
bytes and reserved indices, \code{TX_ABORT}, the abort rule, extra bytes under the same \code{SS},
and finally the transmission and the reception.

It is the \textbf{only} testbench in the course that drives \code{sclk}, \code{mosi} and \code{ss} as
real waveforms. \code{spi_reg_bridge_tb} feeds the bridge finished bytes and never touches a pin, so
everything between the wire and the byte \textemdash{} the synchronisers, the edge detection on SCK,
the \code{ss_active} framing \textemdash{} is exercised here or nowhere.

\subsection{Verification of your own}
\label{app:project:ownverify}

\code{can_spi_node_tb} is deliberately kind to the timing requirements: it holds 200 ns of quiet
around every \code{SS} edge, where the protocol requires at least 60 ns. That is right for a testbench
meant to separate logic errors from timing errors, but it also means that \textbf{no} testbench handed
out presses the specification's timing minima.

That is the verification the project asks the group to do itself: copy \code{can_spi_node_tb}, shrink
the quiet period towards 60 ns, and find out at what value the node stops answering and why. The
answer is in \code{spi_slave}'s synchronisers, and it is the only place where the course lets you
\emph{measure} a margin instead of reading one. Exercise~\ref{c19:ex:node}(b) asks the question in
full.

\subsection{Checklist before hand-in}
\label{app:project:checklist}

\begin{itemize}
  \item All nine testbenches handed out pass, \code{can_controller_tb} and \code{can_spi_node_tb}
    included. \textbf{Check that they really \emph{ran} and were not skipped}, see the box below.
  \item \code{make build-project} goes green from a clean checkout.
  \item The port order matches the tables in \secref{app:project:controller} and
    \ref{app:project:bridge}.
  \item No values hardcoded past \code{can_def}.
  \item Every entity's interface is commented.
  \item \file{CONTRIBUTORS.md} is up to date.
  \item \code{main} contains no commits that did not go through a reviewed PR.
\end{itemize}

\begin{warning}
\textbf{Skipped is not passed.} \code{can_controller_tb} is gated on whether
\file{can_controller.vhd} contains \chapref{c17:ch}'s CRC gating, found by a text search for the
internal signal names the chapter uses. If you write a correct receive path with different names, the
testbench is reported as \emph{skipped}, and since skipped is never failed the checklist looks
satisfied while the course's pass/fail testbench never ran. So run

\begin{shell}
CI_BUILD_ALL=1 make build-project
\end{shell}

at least once before hand-in, and hand in that outcome.
\end{warning}

% ------------------------------------------------------------------------------------------
\section{Examination and grades}
\label{app:project:grading}

The course is examined through \textbf{one group project} and \textbf{two practical exams}. 100 marks
in total.

\begin{booktable}{lllr}
  \thead{Element} & \thead{Where} & \thead{Form} & \thead{Marks} \\
  \midrule
  Practical exam 1, sequential networks & \chapref{c9:ch} & Individual, 3 h & 25 \\
  Group project, a CAN node in VHDL & \chapref{c10:ch}--\ref{c19:ch} & Group of 4--5 & 50 \\
  Practical exam 2, CAN modules & \chapref{c20:ch} & Individual, 3 h & 25 \\
\end{booktable}

\begin{booktable}{lL}
  \thead{Grade} & \thead{Requirement} \\
  \midrule
  \textbf{F} & Below 50 marks, or a failed project. \\
  \textbf{P} & At least 50 marks in total \textbf{and} a passed project. \\
  \textbf{D} & At least 80 marks in total \textbf{and} a distinction-level project. \\
\end{booktable}

The project is therefore a gate in both directions. It accounts for half the course, it is where most
of the work lies, and a failed project cannot be compensated for with two strong exams.

\subsection{The project's 50 marks}
\label{app:project:points}

\begin{booktable}{LrL}
  \thead{Part} & \thead{Marks} & \thead{What is assessed} \\
  \midrule
  The design & 30 & The modules work per the specification and the testbenches pass. \\
  The way of working & 12 & Git history, PRs, the quality of the reviews, the seminar,
    \file{CONTRIBUTORS.md}. \\
  The presentation & 8 & The final presentation: that the group can explain what it built and
    why. \\
\end{booktable}

The grade is given to the group. If the effort is markedly uneven, individual grades may be adjusted;
that is the only reason \file{CONTRIBUTORS.md} is to be kept up to date as you go and not written the
evening before.

\textbf{A passed project} requires a \textbf{complete node}: all nine testbenches handed out must
pass, that is, \code{CI_BUILD_ALL=1 make build-project} green throughout. The controller alone is not
enough, and the reason lies after the course: the node travels on to the CAN labs in the next embedded
systems course, where a driver has to reach it over SPI. A controller without a register bank and a
bridge cannot be reached from any driver at all, and neither can a node whose parts work separately
but not together.

\textbf{Distinction} additionally requires the verification of your own in
\secref{app:project:ownverify}, that the code is documented throughout and stylistically consistent,
and that the review culture is visible in the repository's history.

Bring-up on real hardware is \textbf{not} assessed. A broken kit or a broken connection must not cost
marks. Everything that is assessed can be run in GHDL on an ordinary laptop.

\subsection{The practical exams}
\label{app:project:exams}

Both are individual, three hours, at the computer with GHDL. The form is the same in both:
independent tasks, unrelated to each other and ordered roughly by increasing difficulty.

Most of the tasks are \textbf{writing tasks}: a self-checking testbench handed out and a description
of the module it drives, with port order and behaviour. \Chapref{c20:ch} additionally has
\textbf{reading tasks}: a timing diagram or an SPI transaction to interpret on paper. They have no
testbench, for exactly the reason that makes them worth asking \textemdash{} reading a waveform is a
skill in itself, and it is not examined by writing VHDL. They are marked against an answer key, and
partly correct answers earn partial marks the same way partly correct logic does.

Partly correct logic earns partial marks: never hand in an empty file because the module did not get
finished, and the same goes for a half-interpreted timing diagram. The marking reads the code, not
only the testbench's outcome; a module that passes by doing something other than what was asked does
not earn full marks.

\Chapref{c9:ch} covers \chapref{c1:ch} to \ref{c8:ch}, with the emphasis on sequential networks:
registers, counters, shift registers, timers, synchronisers and simple Moore machines. A practice
paper is handed out in advance. \Chapref{c20:ch} covers the project's modules: bit timing, CRC, bit
stuffing in both directions, frame sequencing, register semantics and the SPI transaction.

\subsection{Permitted aids}
\label{app:project:aids}

This section is the \textbf{authoritative} list. If anything else says otherwise anywhere else, what
is written here applies.

\textbf{Allowed at both exams:}
\begin{itemize}
  \item All the course material: this book, the lectures' READMEs and every appendix. At
    \chapref{c20:ch} the register map and the protocol specification (appendix~\ref{app:protocol})
    belong there of course; at \chapref{c9:ch} they are available too, but are not examined.
  \item All the VHDL you wrote yourself during the course, and at \chapref{c20:ch} the group's
    project code as well.
  \item Your own computer with GHDL, and CircuitVerse. At both exams: drawing a circuit before you
    write it is often faster.
\end{itemize}

\textbf{Not allowed:}
\begin{itemize}
  \item \textbf{Communication with others}, in any form.
  \item \textbf{Tools that generate or explain VHDL for you}: language models, code assistants and
    the like, whether they write the code or merely reason about it. The line is not drawn at who
    presses the keys but at whose understanding is being tested.
  \item \textbf{Searching the internet for solutions.} CircuitVerse runs in a browser, so the computer
    is online, and that assumes you keep to the course material under your own steam. Looking up
    GHDL's or VHDL's documentation is allowed; searching for a finished module is not.
\end{itemize}

If you are unsure whether something is allowed: ask before the exam, not during it.

That the group's project code is allowed at \chapref{c20:ch} is a deliberate choice. You built it. The
tasks are instead written as variants and repairs, so that they differ from the project's modules at
precisely the point that requires you to have understood them.

\subsection{Re-examination and absence}
\label{app:project:reexam}

Both exams can be retaken. The project is handed in once, with the possibility of a resubmission
after word from the supervisor: a resubmission is marked against a pass, not against a distinction.

The code review seminar in \chapref{c15:ch} and the final presentation in \chapref{c19:ch} are the
project's two compulsory checkpoints. Missed attendance is made up by agreement with the supervisor.
Other sessions have no attendance requirement; the group project, on the other hand, has a group
dynamic that does not survive somebody being absent, which is a different and stronger reason to come.

% ------------------------------------------------------------------------------------------
\section{Practical tips}
\label{app:project:tips}

\begin{itemize}
  \item Draw the timing diagram before you write the shift registers. Most faults there are one-cycle
    timing errors, not logic errors, and they show up on paper before they show up in the simulator.
  \item Start with \code{bit_timer} and \code{crc15} if you are unsure where to start. They are the
    simplest, they depend on nothing, and they can be finished completely.
  \item Read the testbenches handed out. They say more precisely what every module promises than any
    prose does, and they are written to be read.
  \item Test the stuffing carefully, in both directions. It is the project's most common source of
    failure.
  \item Keep \code{can_spi_node} structural. As soon as logic creeps into the top level it becomes
    impossible to test in isolation.
  \item Talk to the parallel class early, and do it about the register map rather than about your
    code.
\end{itemize}

\begin{aside}
\textbf{Shared documentation.} Two documents are shared with the software group and are the contract
between the classes: the register map (what the registers mean) and the protocol specification (how
they are reached), both reproduced in appendix~\ref{app:protocol}. If anything in them changes, both
classes have to know about it before the change is merged. The documents are updated \emph{before} the
code, never after.
\end{aside}
